\section{Grafi: Visite DFS e BFS (Lab 21)}

\subsection{Attenzione ai Cicli!}
A differenza degli alberi, nei grafi puoi girare in tondo.
È \textbf{obbligatorio} usare un insieme \texttt{visited} per ricordarsi dove siamo già stati, altrimenti il programma va in loop infinito.

\subsection{DFS (Visita in Profondità)}
\begin{concept}{Labirinto (Filo di Arianna)}
Esplora un sentiero fino alla fine, poi torna indietro (**Backtracking**).
Usa la **Ricorsione** (Stack).
\end{concept}
\begin{jsbox}
function dfs(nodo, visited = new Set()) {
    // 1. Segno come visitato
    visited.add(nodo);
    console.log(nodo);
    
    // 2. Prendo i vicini dalla Mappa (vedi Cap. 20)
    let vicini = grafo.get(nodo);
    
    // 3. Per ogni vicino non ancora visto, vado in profondità
    for (let vicino of vicini) {
        if (!visited.has(vicino)) {
            dfs(vicino, visited);
        }
    }
}
\end{jsbox}

\subsection{BFS (Visita in Ampiezza)}
\begin{concept}{Onda d'Urto}
Esplora a cerchi concentrici (tutti i vicini a distanza 1, poi distanza 2...).
Usa una **Coda** (FIFO).
\end{concept}
\begin{jsbox}
function bfs(startNode) {
    let visited = new Set();
    let coda = [startNode];
    
    visited.add(startNode); // Importante: segna subito!
    
    while (coda.length > 0) {
        let nodo = coda.shift(); // Prendo il primo
        console.log(nodo);
        
        let vicini = grafo.get(nodo);
        for (let v of vicini) {
            if (!visited.has(v)) {
                visited.add(v); // Segna
                coda.push(v);   // Aggiungi alla coda
            }
        }
    }
}
\end{jsbox}

\subsection{Applicazioni}
\begin{itemize}
    \item **Componenti Connesse:** Se lanci una DFS e non visiti tutti i nodi, significa che il grafo è spezzato in isole.
    \item **Shortest Path:** La BFS trova sempre la strada col minor numero di salti.
\end{itemize}
